\documentclass[a4paper,11pt]{article}
\usepackage[utf8]{inputenc}
\usepackage[T1]{fontenc}
\usepackage{lmodern}
\usepackage{amsmath,amssymb,amsthm}
\usepackage{mathtools}
\usepackage{booktabs}
\usepackage{longtable}
\usepackage{array}
\usepackage{hyperref}
\usepackage{graphicx}
\usepackage{float}
\usepackage{geometry}
\geometry{margin=25mm}

\hypersetup{
  colorlinks=true,
  linkcolor=blue,
  urlcolor=blue,
  citecolor=blue
}

\theoremstyle{definition}
\newtheorem{theorem}{Theorem}[section]
\newtheorem{proposition}[theorem]{Proposition}
\newtheorem{lemma}[theorem]{Lemma}
\newtheorem{corollary}[theorem]{Corollary}
\newtheorem{definition}[theorem]{Definition}
\newtheorem{remark}[theorem]{Remark}
\newtheorem{observation}[theorem]{Observation}

\setlength{\parindent}{1em}
\setlength{\parskip}{0.5em}

\providecommand{\tightlist}{\setlength{\itemsep}{0pt}\setlength{\parskip}{0pt}}
\begin{document}

\section{On the Continuity of Geodesics in Subjective Spaces}\label{on-the-continuity-of-geodesics-in-subjective-spaces}

\textbf{--- On Bridging Geodesics from a Parent Subjective Space to a Child Subjective Space ---}

\begin{center}\rule{0.5\linewidth}{0.5pt}\end{center}

\textbf{Author}: Noriaki Kihara (WF System Co., Ltd.)

\textbf{Date}: April 2026

\textbf{DOI:} \href{https://doi.org/10.5281/zenodo.19533299}{10.5281/zenodo.19533299}

\begin{center}\rule{0.5\linewidth}{0.5pt}\end{center}

\subsection{Abstract}\label{abstract}

Based on the central projection framework established in prior papers {[}1{]}, {[}2{]}, {[}3{]}, {[}4{]}, {[}5{]}, {[}6{]}, this paper examines the behavior of geodesics on the subjective space as they are traced in the direction \(R \to 0\), and discusses the bridging of geodesics within the nested structure demonstrated in Proposition 4.1 of paper {[}4{]}.

\textbf{Disclaimer}: This paper proves nothing new. The definitions and properties of geodesics on the subjective space have already been organized in papers {[}1{]} and {[}2{]}, and the existence of the nested structure has already been established in Proposition 4.1 of paper {[}4{]}. This paper merely provides a geometric description of the relationship between geodesics in the parent subjective space and geodesics in the child subjective space, as a combination of these existing results. It only organizes the correspondence with the established notion of geodesics in differential geometry, and the correspondence with concepts in physics such as geodesic completeness, singularity theorems, and black hole singularities is outside the scope of this paper.

\begin{center}\rule{0.5\linewidth}{0.5pt}\end{center}

\subsection{1. Introduction}\label{introduction}

In the central projection framework introduced in prior paper {[}1{]}, the subjective space \(\Pi_R\) was defined as a space governed by the curvature radius \(R\). Paper {[}2{]} organized the properties of geodesics (great circles) on the subjective space and the geodesic deviation \(|\xi|^2\), and paper {[}4{]} demonstrated geometric regularity at both limits \(R \to 0\) and \(R \to \infty\), and proved the existence of a nested structure of subjective spaces in Proposition 4.1. Paper {[}6{]} organized the geometric properties of subjective spaces of each dimension as correspondences with established results.

The purpose of this paper is to combine the results of paper {[}4{]} with the geodesic properties of paper {[}2{]} in order to organize how, when a geodesic on the parent subjective space \(\Pi_R\) is traced in the direction \(R \to 0\), the bridging to the child subjective space \(\Pi_{R'}\) guaranteed by Proposition 4.1 of paper {[}4{]} is described.

This paper proves nothing new. It only provides a description that combines Proposition 4.1 of paper {[}4{]} with the geodesic properties of paper {[}2{]}.

\begin{center}\rule{0.5\linewidth}{0.5pt}\end{center}

\subsection{2. Preliminaries: Geodesics on the Subjective Space}\label{preliminaries-geodesics-on-the-subjective-space}

We follow the terminology of paper {[}1{]} §2 and paper {[}2{]} §2. The principal terms used in this paper are reproduced below.

{\def\LTcaptype{none} % do not increment counter
\begin{longtable}[]{@{}
  >{\raggedright\arraybackslash}p{(\linewidth - 4\tabcolsep) * \real{0.3333}}
  >{\raggedright\arraybackslash}p{(\linewidth - 4\tabcolsep) * \real{0.3333}}
  >{\raggedright\arraybackslash}p{(\linewidth - 4\tabcolsep) * \real{0.3333}}@{}}
\toprule\noalign{}
\begin{minipage}[b]{\linewidth}\raggedright
Term
\end{minipage} & \begin{minipage}[b]{\linewidth}\raggedright
Definition
\end{minipage} & \begin{minipage}[b]{\linewidth}\raggedright
Symbol
\end{minipage} \\
\midrule\noalign{}
\endhead
\bottomrule\noalign{}
\endlastfoot
Subjective space & The image of the central projection \(\Phi_A\). The space inhabited by the internal observer. & \(\Pi_R\) or \(H^+_A\) \\
Curvature radius & The radius governing the subjective space. & \(R\) \\
Geodesic & A great circle on the subjective space. Corresponds to a straight line on the tangent hyperplane by Proposition 2.1 of paper {[}1{]}. & --- \\
Geodesic deviation & The squared norm of the deviation vector \(\xi^\mu\) between adjacent geodesics. & \(\vert\xi\vert^2 = g_{\mu\nu}\xi^\mu\xi^\nu\) \\
\end{longtable}
}

By Proposition 2.1 of paper {[}1{]}, under the central projection, a straight line on the tangent hyperplane maps to a great circle (geodesic) on the subjective space \(\Pi_R\). By Theorem 5.1 and Proposition 5.1 of paper {[}2{]}, the geodesic deviation satisfies the Jacobi equation, and \(|\xi|^2\) is a scalar invariant of dimension {[}L²{]} measurable by the internal observer of the subjective space.

By Theorem 5.2 (i) of paper {[}2{]}, there exists an upper bound on the measured value of \(|\xi|^2\) that depends on the curvature radius \(R\). This upper bound is hereafter denoted \(|\xi|^2_{\max}(R)\).

\begin{center}\rule{0.5\linewidth}{0.5pt}\end{center}

\subsection{3. Tracing Geodesics in the Parent Subjective Space}\label{tracing-geodesics-in-the-parent-subjective-space}

\subsubsection{3.1 Tracing a Geodesic and Geodesic Deviation}\label{tracing-a-geodesic-and-geodesic-deviation}

Fix a geodesic \(\gamma\) on the parent subjective space \(\Pi_R\). By Proposition 2.1 of paper {[}1{]}, \(\gamma\) corresponds to a straight line on the tangent hyperplane \(\Pi^{(A)}_R\) and is globally well-defined. Moreover, by Theorem 3.1 of paper {[}2{]} (discrete stability), every point on \(\gamma\) lies in the regular domain of the map.

The deviation \(\xi^\mu\) between \(\gamma\) and adjacent geodesics behaves according to the Jacobi equation of Proposition 5.1 of paper {[}2{]}. The measured value of \(|\xi|^2\) is constrained by the upper bound \(|\xi|^2_{\max}(R)\) of Theorem 5.2 (i) of paper {[}2{]}.

\subsubsection{\texorpdfstring{3.2 The Direction \(R \to 0\) and the Degeneration of the Measurement Upper Bound}{3.2 The Direction R \textbackslash to 0 and the Degeneration of the Measurement Upper Bound}}\label{the-direction-r-to-0-and-the-degeneration-of-the-measurement-upper-bound}

In paper {[}4{]} §3, it was noted that in the limit \(R \to 0\), the upper bound \(|\xi|^2_{\max}(R)\) on the geodesic deviation asymptotes to zero, while the local geometric structure of the subjective space itself remains well-defined for \(R > 0\) --- a structural tension. This tension arises between the quantities measurable by an observer inside the subjective space and the intrinsic structure of the subjective space that observer inhabits.

The observation in this section is a reproduction of paper {[}4{]} §3.2 and contains no new claims.

\subsubsection{3.3 The Status of the Geodesic as Seen from Inside the Parent Subjective Space}\label{the-status-of-the-geodesic-as-seen-from-inside-the-parent-subjective-space}

Following the definition of the internal observer in paper {[}2{]} §2, the information available to the internal observer of the parent subjective space \(\Pi_R\) is restricted to scalar invariants constructed from the metric \(g_{\mu\nu}\) on \(\Pi_R\) (Proposition 3.1 of paper {[}3{]}). The \(|\xi|^2\) measured by the internal observer along \(\gamma\) is well-defined within the range not exceeding \(|\xi|^2_{\max}(R)\).

The status of \(\gamma\) in the degeneration region of \(|\xi|^2_{\max}(R)\) falls within the range of the structural tension of paper {[}4{]} §3.2, and cannot be fully resolved within a framework that considers only the interior of the parent subjective space.

\begin{center}\rule{0.5\linewidth}{0.5pt}\end{center}

\subsection{4. Bridging to the Nested Structure}\label{bridging-to-the-nested-structure}

\subsubsection{4.1 Reproduction of Proposition 4.1 from Paper {[}4{]}}\label{reproduction-of-proposition-4.1-from-paper-4}

Proposition 4.1 of paper {[}4{]} (existence of the nested structure) is reproduced here in the context of this paper.

\textbf{Proposition 4.1 (from paper {[}4{]})} In the central projection framework defined in paper {[}1{]}, for any subjective space \(\Pi_R\) and any point \(p \in \Pi_R\) in its interior, a subjective space \(\Pi_{R'}\) satisfying the following properties exists for any \(R' > 0\):

\begin{enumerate}
\def\labelenumi{\arabic{enumi}.}
\tightlist
\item
  \(\Pi_{R'}\) is defined by a new central projection centered at \(p\).
\item
  \(\Pi_{R'}\) has the same geometric structure as the construction defined in paper {[}1{]}.
\item
  \(\Pi_{R'}\) is constructed from the same background space \(\mathbb{R}^{n+1}\) as \(\Pi_R\).
\end{enumerate}

\subsubsection{4.2 Points on the Geodesic as Seed Points}\label{points-on-the-geodesic-as-seed-points}

For any point \(p\) on the geodesic \(\gamma\) traced in §3, Proposition 4.1 guarantees that it is geometrically possible to instantiate a child subjective space \(\Pi_{R'}\) centered at \(p\). The instantiated child subjective space \(\Pi_{R'}\) carries the same pullback metric, curvature tensor, and Einstein tensor as the construction of paper {[}1{]}.

Within the interior of the child subjective space \(\Pi_{R'}\), the results of papers {[}1{]} and {[}2{]} apply directly. Geodesics on \(\Pi_{R'}\) are defined as great circles according to Proposition 2.1 of paper {[}1{]}, and their geodesic deviation follows the Jacobi equation of Proposition 5.1 of paper {[}2{]}. The \(|\xi|^2\) measurable by the internal observer of \(\Pi_{R'}\) is constrained by the upper bound \(|\xi|^2_{\max}(R')\) that depends on \(R'\).

\subsubsection{4.3 The Relationship Between the Parent Geodesic and the Child Geodesic}\label{the-relationship-between-the-parent-geodesic-and-the-child-geodesic}

By Symmetry IV of paper {[}2{]} (transformability of subjective coordinate systems, Theorem 6.1), the subjective coordinates of the parent subjective space \(\Pi_R\) and those of the child subjective space \(\Pi_{R'}\) are related by a regular composite map passing through a common point in the background space \(\mathbb{R}^{n+1}\).

In the child subjective space \(\Pi_{R'}\) seeded at point \(p\) on the geodesic \(\gamma\) of the parent subjective space \(\Pi_R\), a geodesic \(\gamma'\) departing from \(p\) is well-defined on \(\Pi_{R'}\) according to the results of papers {[}1{]} and {[}2{]}. The geodesics \(\gamma\) and \(\gamma'\) share \(p\) as a common point and are independently defined geodesics within their respective subjective spaces.

The behavior of \(\gamma\) as seen by the internal observer of the parent subjective space \(\Pi_R\), and the behavior of \(\gamma'\) as seen by the internal observer of the child subjective space \(\Pi_{R'}\), are each described as quantities constructed from the metric within their respective subjective spaces, by Theorem 5.1 of paper {[}3{]} (relativity of observation). Although the two share the same background space, they are each independently defined geometric objects within their respective subjective spaces.

\subsubsection{4.4 Scope of the Claims in This Section}\label{scope-of-the-claims-in-this-section}

What has been stated in this section is limited to the following.

\begin{enumerate}
\def\labelenumi{\arabic{enumi}.}
\tightlist
\item
  That for any point \(p\) on the geodesic \(\gamma\) of the parent subjective space \(\Pi_R\), taken as the seed point, the construction of the child subjective space \(\Pi_{R'}\) is geometrically possible (a direct consequence of Proposition 4.1 of paper {[}4{]}).
\item
  That the geodesic \(\gamma'\) is well-defined on the child subjective space \(\Pi_{R'}\) as well, according to the results of papers {[}1{]} and {[}2{]} (a direct consequence of Proposition 2.1 of paper {[}1{]} and Proposition 5.1 of paper {[}2{]}).
\item
  That the parent geodesic \(\gamma\) and the child geodesic \(\gamma'\), while sharing the same background space, are independently defined geometric objects within their respective subjective spaces (a direct consequence of Theorem 5.1 of paper {[}3{]}).
\end{enumerate}

This section contains no new proofs. These claims are described as a combination of existing propositions from papers {[}1{]}--{[}4{]}.

\begin{center}\rule{0.5\linewidth}{0.5pt}\end{center}

\subsection{5. Discussion}\label{discussion}

In this paper, we have organized the behavior of the geodesic \(\gamma\) on the parent subjective space \(\Pi_R\) when traced in the direction \(R \to 0\), and the geometric relationship of the instantiation into a child subjective space guaranteed by Proposition 4.1 of paper {[}4{]}.

What has been described in this paper is limited to the following.

\begin{enumerate}
\def\labelenumi{\arabic{enumi}.}
\tightlist
\item
  That in tracing geodesics on the parent subjective space \(\Pi_R\), the upper bound on geodesic deviation from Theorem 5.2 (i) of paper {[}2{]} and the structural tension of paper {[}4{]} §3.2 emerge.
\item
  That by Proposition 4.1 of paper {[}4{]}, the construction of a child subjective space \(\Pi_{R'}\) seeded at any point on the geodesic of the parent subjective space is geometrically possible.
\item
  That geodesics are well-defined on the child subjective space as well according to the results of papers {[}1{]} and {[}2{]}, and that the parent and child geodesics, while sharing the same background space, are each independently defined within their respective subjective spaces.
\end{enumerate}

Each item described in this paper is a direct combination of existing propositions from papers {[}1{]}--{[}4{]} and contains no new proofs. The intent of this paper is limited to complementing the description of the nested structure of paper {[}4{]} by focusing on geodesics as geometric objects.

The physical significance these observations might carry is not addressed in this paper. In particular, the correspondence with the established notion of geodesic completeness in differential geometry, the singularity theorems in general relativity, and the behavior of geodesics in black hole spacetimes are matters outside the geometric framework and beyond the scope of this paper. Answers to these questions are left to the judgment of the reader.

\begin{center}\rule{0.5\linewidth}{0.5pt}\end{center}

\subsection{6. Conclusion}\label{conclusion}

Geodesics on the parent subjective space \(\Pi_R\) of the central projection framework are well-defined according to the results of paper {[}2{]} and carry the structural tension of paper {[}4{]} §3.2. By Proposition 4.1 of paper {[}4{]}, a child subjective space \(\Pi_{R'}\) seeded at any point on the geodesic of the parent subjective space can be constructed, and geodesics are independently well-defined on the child subjective space as well. The parent geodesic and the child geodesic, while sharing the same background space, are independently defined geometric objects within their respective subjective spaces. All of these are described as direct combinations of existing propositions from papers {[}1{]}--{[}4{]}. Physical interpretation is outside the scope of this paper and is left to the reader.

\begin{center}\rule{0.5\linewidth}{0.5pt}\end{center}

\subsection{References}\label{references}

{[}1{]} Noriaki Kihara, ``Geometric Formulation of Four-Dimensional Space via Central Projection'', Zenodo, DOI: 10.5281/zenodo.19427780 (2026).

{[}2{]} Noriaki Kihara, ``Geometric Symmetries of Central Projection: Mathematical Foundations of the Multi-Axis Model'', Zenodo, DOI: 10.5281/zenodo.19434932 (2026).

{[}3{]} Noriaki Kihara, ``Relativity of Observation in Multiple Subjective Spaces: Geometric Consequences of the Symmetry of Central Projection'', Zenodo, DOI: 10.5281/zenodo.19435162 (2026).

{[}4{]} Noriaki Kihara, ``On the Limits of the Curvature Radius in Central Projection --- The Cases \(R \to \infty\) and \(R \to 0\)'', Zenodo, DOI: 10.5281/zenodo.19526549 (2026).

{[}5{]} Noriaki Kihara, ``On Dimension Addition to Background Spaces and Subjective Spaces'', Zenodo, DOI: 10.5281/zenodo.19526913 (2026).

{[}6{]} Noriaki Kihara, ``Subjective Spaces from Zero to Four Dimensions'', Zenodo, DOI: 10.5281/zenodo.19533292 (2026).

\begin{center}\rule{0.5\linewidth}{0.5pt}\end{center}

\end{document}
